\chapter*{Identificación del documento}
\addcontentsline{toc}{chapter}{Identificación del documento}

\begin{table}[H]
\caption{Ficha de identificación del informe}
\label{tab:identificacion}
\centering
\begingroup
\rowcolors{2}{SoftGray}{white}
\begin{tabularx}{\textwidth}{>{\bfseries\color{EspeGreenDark}}p{0.30\textwidth}X}
\toprule
Documento & Informe técnico final de InvenTrack \\
Autor & \ReportAuthor \\
Asignatura & Desarrollo Seguro \\
Institución & Universidad de las Fuerzas Armadas ESPE \\
Período académico & \AcademicPeriod \\
Fecha del informe & \ReportDate \\
Entorno & WSL Ubuntu, Kali Linux WSL, Docker, Minikube y Kubernetes \\
Dominio del laboratorio & \url{http://conjunta3p.espe.edu.ec} \\
Repositorio del informe & \href{\ReportRepo}{\ReportRepo} \\
Repositorio Kubernetes & \href{\KubernetesRepo}{\KubernetesRepo} \\
Repositorio de la aplicación & \href{\ApplicationRepo}{\ApplicationRepo} \\
\bottomrule
\end{tabularx}
\endgroup
\end{table}

\vspace{0.45cm}

\begin{tcolorbox}[infobox]
\textbf{Trazabilidad.} El cuerpo del informe resume los resultados y conserva
los comandos ejecutados. El Anexo A funciona como mapa de rutas hacia cada
archivo de evidencia incluido en el repositorio del informe.
\end{tcolorbox}

\clearpage
\tableofcontents
\clearpage
\listoftables
\clearpage
